
Benchmark contamination---the presence of evaluation data in training corpora---threatens the reliability of LLM evaluation. While contamination detection methods exist, it remains unclear how data curation decisions affect contamination rates. This work investigates whether perplexity-based filtering, a standard technique in modern data curation, systematically amplifies benchmark contamination. We introduce Contamination Contribution Ratio (CCR) to quantify how much benchmark performance derives from contaminated training examples, combining n-gram contamination detection with data attribution. In methodology validation experiments, synthetic injection produces linear CCR scaling ($R^2$ = 0.9998), demonstrating CCR as a calibrated metric. Simulated filtering comparisons show perplexity filtering increases CCR by 0.1594 relative to random sampling (p < 0.0001). Removal intervention experiments indicate that high-CCR examples are causally necessary for benchmark performance: removing the top 1-5\% of high-CCR examples produces 1.$97\times$ greater accuracy degradation than random removal (95\% CI: [1.53, 2.34]). Additionally, contaminated examples exhibit 4.$1\times$ higher Influence Fragility Ratio than non-contaminated examples, suggesting structural irreplaceability. However, the hypothesized negative correlation between influence fragility and k-NN redundancy was weaker than expected ($\rho$ = -0.11 vs. threshold $\rho$ < -0.5). All results reported here derive from methodology validation using simulated attribution due to GPU infrastructure constraints; full empirical validation with model training is ongoing.
